\chapter{Điều hướng bối cảnh quy định và đạo đức của AI trong học thuật}

\subsection{Tóm tắt}

Chương 8 xem xét những phức tạp về quy định pháp lý và đạo đức xung quanh việc sử dụng AI tạo sinh trong môi trường học thuật. Khi các công cụ AI ngày càng ảnh hưởng đến nghiên cứu, giảng dạy và đánh giá, giới học giả phải xoay xở trong một bối cảnh pháp lý phân mảnh và biến đổi nhanh chóng, bao gồm các vấn đề về bảo vệ dữ liệu, sở hữu trí tuệ và tuân thủ quy định của tổ chức. Chương này làm nổi bật các rủi ro pháp lý chính, chẳng hạn như vi phạm quyền riêng tư theo GDPR và những bất định về bản quyền gắn với nội dung do AI tạo ra. Chương cũng đề cập đến các mối quan ngại đạo đức sâu xa hơn, gồm xu hướng của AI làm gia tăng thiên kiến học thuật, tính thiếu minh bạch của nó, cùng tác động đối với quyền tác giả, tính độc đáo và việc đánh giá sinh viên. Dựa trên các ví dụ từ phản biện đồng cấp (peer review), hồ sơ xin tài trợ và môi trường giáo dục, chương này nhấn mạnh sự cần thiết của các chính sách minh bạch, sự tiếp cận mang tính phản biện và định hướng từ tổ chức nhằm bảo đảm AI tạo sinh được sử dụng một cách có trách nhiệm. Thay vì đưa ra những câu trả lời dứt khoát, chương trang bị cho giới học thuật các công cụ khái niệm và thực tiễn cần thiết để cân bằng giữa đổi mới và liêm chính học thuật.

\subsection{Từ khóa}

quy định về AI, liêm chính học thuật, đạo đức của AI tạo sinh, bản quyền, bảo vệ dữ liệu, GDPR, quyền tác giả, phản biện đồng cấp (peer review), đánh giá sinh viên, tuân thủ của tổ chức, xuất bản học thuật, thiên kiến của AI, tính minh bạch của AI, sở hữu trí tuệ, chính sách học thuật

Việc tích hợp AI tạo sinh vào môi trường học thuật đặt ra những thách thức pháp lý và đạo đức sâu sắc mà các nhà nghiên cứu, nhà giáo dục và các tổ chức hiện phải đối mặt. Những thách thức này đặc biệt phức tạp vì công nghệ AI phát triển nhanh chóng, trong khi các khung pháp lý và chính sách của tổ chức khó theo kịp. Các trường đại học và nhà xuất bản đang xây dựng hướng dẫn, tòa án đang đưa ra những phán quyết mới, và các nhà hoạch định chính sách đang tranh luận về cách cân bằng giữa lợi ích và rủi ro của AI (Jabotinsky \& Sarel, 2023). Trong khi đó, các học giả ở mọi giai đoạn sự nghiệp phải xoay xở trong một bối cảnh quy tắc và kỳ vọng luôn thay đổi, thường không có hướng dẫn rõ ràng về điều gì được phép, nên làm hay phù hợp về mặt đạo đức.

Sự bất định này làm nảy sinh hàng loạt câu hỏi cấp bách: Các nhà nghiên cứu có được phép sử dụng AI trong công việc của mình không, và nếu có thì những công cụ nào được chấp nhận? Nội dung do AI tạo ra có thể được đưa vào các ấn phẩm không, và có nên công khai việc sử dụng đó không? Nếu các mô hình AI đã được huấn luyện trên tài liệu học thuật có bản quyền, điều này có ảnh hưởng đến cách các nhà nghiên cứu nên tương tác với chúng không (Lemley, 2023)? Làm thế nào các học giả có thể bảo đảm tính chính xác và tính liêm chính của nội dung do AI tạo ra? Hệ quả đối với việc học tập và đánh giá của sinh viên là gì khi AI có thể viết tiểu luận, phân tích dữ liệu và giải quyết các vấn đề phức tạp? Các tổ chức nên định nghĩa việc sử dụng AI có trách nhiệm trong nghiên cứu, giảng dạy và công tác hành chính như thế nào? Hơn nữa, các mối lo ngại về quyền riêng tư là rất lớn, đặc biệt khi dữ liệu nghiên cứu nhạy cảm được xử lý thông qua các hệ thống AI bên ngoài. Ở tầm rộng hơn, các câu hỏi về tính xác thực học thuật xuất hiện trong thời đại mà nội dung do AI tạo ra có thể làm mờ ranh giới của nỗ lực trí tuệ của con người (U.K. Intellectual Property Office, 2024).

Bối cảnh quy định xung quanh AI tạo sinh mang tính năng động và phân mảnh. Các khu vực pháp lý khác nhau áp dụng những cách tiếp cận khác nhau đối với quản trị AI, và chính sách của các tổ chức tiếp tục thay đổi khi các trường đại học vật lộn tìm kiếm thực tiễn tốt nhất (European Commission, 2023). Tuy nhiên, trước khi khám phá các thách thức quy định chính, cần nhấn mạnh rằng chương này không cấu thành lời tư vấn pháp lý. Các nhà nghiên cứu và sinh viên luôn nên tham khảo chính sách của tổ chức mình, cố vấn pháp lý hoặc các cơ quan tài trợ để có hướng dẫn cụ thể phù hợp với hoàn cảnh của mình. Thay vào đó, phần thảo luận này cung cấp một cái nhìn tổng quan về các cân nhắc quy định và đạo đức then chốt đang định hình việc sử dụng AI trong môi trường học thuật. Với tốc độ thay đổi trong lĩnh vực này, những gì hiện được phép có thể bị hạn chế vào ngày mai, và các tổ chức có thể áp dụng những yêu cầu mới khi năng lực của AI mở rộng.

Tuy nhiên, không phải mọi mối quan ngại xung quanh AI đều thuần túy mang tính pháp lý. Các vấn đề nan giải về đạo đức do AI tạo sinh đặt ra vượt xa việc tuân thủ luật bảo vệ dữ liệu hay các hạn chế về bản quyền. Những công cụ này có khả năng củng cố các thiên kiến vốn có trong dữ liệu huấn luyện, ảnh hưởng đến kết luận nghiên cứu hoặc gạt ra bên lề các quan điểm ít được đại diện (Bommasani et al., 2021). Nội dung do AI tạo ra có thể đưa vào các lỗi hoặc thông tin gây hiểu lầm, có nguy cơ làm lan truyền những điều không chính xác trong diễn ngôn học thuật (Cotton et al., 2023). Trong dạy và học, sự trỗi dậy của các công cụ viết có hỗ trợ của AI đặt ra những câu hỏi căn bản về tính công bằng của đánh giá, kỹ năng tư duy phản biện và vai trò của sự sáng tạo của con người trong công việc học thuật (Dehouche, 2021). Các vấn đề đạo đức cũng mở rộng sang phương pháp luận nghiên cứu: nếu AI có thể tạo ra các bài tổng quan tài liệu hoặc tổng hợp lập luận, liệu điều này có chuẩn hóa các cách tiếp cận nghiên cứu theo hướng làm nản lòng tính độc đáo về trí tuệ không? Ngoài ra, các vấn đề về quyền tác giả, ghi nhận nguồn đóng góp và tính minh bạch đang trở thành trung tâm khi vai trò của AI trong công việc học thuật ngày càng mở rộng (Jabotinsky \& Sarel, 2023).

Chương này khám phá cả khía cạnh pháp lý và đạo đức của AI tạo sinh trong môi trường học thuật. Dù không thể đưa ra câu trả lời dứt khoát cho tất cả những thách thức đang tiến triển này, chương nêu bật các cân nhắc then chốt mà các học giả và tổ chức phải đối mặt. Thay vì cố gắng dự đoán mọi thay đổi về quy định hay chính sách của tổ chức, phần thảo luận này tập trung vào những vấn đề cấp bách nhất đang định hình vị trí của AI trong nghiên cứu, giảng dạy và xuất bản học thuật ngày nay. Hiểu rõ những vấn đề này là điều thiết yếu đối với bất kỳ ai muốn sử dụng các công cụ AI mà vẫn duy trì liêm chính học thuật, tuân thủ pháp luật và trách nhiệm đạo đức.

\section{Những thách thức về quy định pháp lý}

Việc sử dụng AI tạo sinh trong môi trường học thuật làm nảy sinh những vấn đề pháp lý và quy định phức tạp. Các luật bảo vệ dữ liệu, chẳng hạn Quy định chung về bảo vệ dữ liệu (General Data Protection Regulation, GDPR) tại Liên minh châu Âu, đặt ra những yêu cầu nghiêm ngặt khi các công cụ AI xử lý dữ liệu cá nhân. Các mối quan ngại về sở hữu trí tuệ làm phức tạp thêm những câu hỏi về quyền sở hữu, bản quyền và tính hợp pháp của các tập dữ liệu huấn luyện AI. Các chính sách liêm chính học thuật ngày càng yêu cầu minh bạch trong nghiên cứu và viết có sự hỗ trợ của AI. Ngoài ra, các nghĩa vụ hợp đồng với nhà xuất bản và những hạn chế về cấp phép đối với các công cụ AI cũng định hình cách thức các công nghệ này có thể được sử dụng trong bối cảnh học thuật. Khi các tổ chức và cơ quan quản lý thích ứng với vai trò ngày càng lớn của AI, giới học giả phải định hướng trong một bối cảnh các ràng buộc pháp lý và đạo đức không ngừng thay đổi.

\subsection{\texorpdfstring{\textbf{Bảo vệ dữ liệu và quyền riêng tư}}{Bảo vệ dữ liệu và quyền riêng tư}}

Việc sử dụng AI tạo sinh trong nghiên cứu học thuật làm dấy lên các mối quan ngại về quyền riêng tư và bảo vệ dữ liệu, đặc biệt khi nhiều công cụ AI xử lý dữ liệu đầu vào của người dùng trên các máy chủ bên ngoài. Điều này có nghĩa là bất cứ thứ gì được nhập vào một hệ thống AI đều có thể bị lưu trữ, phân tích, thậm chí được dùng để huấn luyện mô hình trong tương lai, trừ khi có tuyên bố rõ ràng khác (Ye et al., 2024). Các nhà nghiên cứu xử lý thông tin nhạy cảm, bao gồm dữ liệu cá nhân, bản thảo mang tính bảo mật hoặc các phát hiện chưa công bố, phải thực hiện thêm các biện pháp phòng ngừa để bảo đảm tuân thủ các quy định về bảo vệ dữ liệu và chính sách của tổ chức. Nếu thiếu các biện pháp bảo vệ thích đáng, thông tin nhạy cảm có thể vô tình bị lộ, bị lạm dụng hoặc bị các bên không được phép truy cập (European Data Protection Supervisor, 2023).

Một trong những khuôn khổ pháp lý chính điều chỉnh việc bảo vệ dữ liệu là GDPR, được áp dụng tại EU và có ảnh hưởng đến các chuẩn mực về quyền riêng tư dữ liệu trên toàn thế giới. Theo GDPR, mọi hoạt động xử lý dữ liệu cá nhân phải được thực hiện một cách hợp pháp, công bằng và minh bạch (van Ooijen \& Vrabec, 2019). Nếu các nhà nghiên cứu sử dụng công cụ AI để xử lý thông tin có thể nhận dạng được, họ phải bảo đảm có cơ sở pháp lý để làm như vậy, đặc biệt khi thực hiện nghiên cứu trên đối tượng là con người. Điều này không chỉ áp dụng cho dữ liệu của người tham gia mà còn cho các nghiên cứu của tổ chức có thể chịu ràng buộc bởi các thỏa thuận bảo mật. Việc chuyển dữ liệu cá nhân có thể nhận dạng được sang các nền tảng AI bên ngoài mà không có sự đồng ý rõ ràng có thể cấu thành vi phạm luật bảo vệ dữ liệu, ngay cả khi không xảy ra việc lạm dụng trực tiếp.

Để giảm thiểu các rủi ro này, các nhà nghiên cứu nên sử dụng những mô hình AI cam kết rõ ràng về việc bảo vệ quyền riêng tư. Nhiều công cụ AI đa dụng lưu trữ dữ liệu đầu vào của người dùng theo mặc định, nhất là ở các phiên bản miễn phí. Chẳng hạn, phiên bản miễn phí của ChatGPT lưu giữ các cuộc hội thoại và có thể dùng chúng để huấn luyện các mô hình trong tương lai, trong khi ChatGPT Enterprise không lưu trữ dữ liệu đầu vào của người dùng. Tương tự, một số công cụ AI học thuật đưa ra cam kết rõ ràng rằng dữ liệu sẽ không được lưu giữ. Trước khi sử dụng bất kỳ công cụ AI nào, các nhà nghiên cứu nên xem lại cài đặt quyền riêng tư của nền tảng và, khi cần thiết, hỏi trực tiếp xem mô hình có lưu giữ dữ liệu hoặc dùng dữ liệu để huấn luyện hay không. Một số nhà cung cấp AI cho phép người dùng tắt tính năng lưu trữ dữ liệu hoặc lịch sử, nhưng các cài đặt này thường không được bật theo mặc định.

Đối với nghiên cứu liên quan đến dữ liệu nhạy cảm hoặc chưa công bố, ẩn danh hóa có thể đóng vai trò như một biện pháp bảo vệ, nhưng đây không phải là giải pháp tuyệt đối an toàn. Thay vì nhập toàn bộ mô tả nghiên cứu hoặc thông tin chi tiết về người tham gia, có thể dùng các ký hiệu giữ chỗ hoặc bút danh để loại bỏ thông tin nhận dạng. Tuy nhiên, các mô hình AI có thể suy luận những chi tiết bị thiếu dựa trên ngữ cảnh, và khi đối chiếu chéo với các nguồn khác, ngay cả dữ liệu mới được ẩn danh một phần đôi khi vẫn có thể bị tái nhận dạng (Wachter \& Mittelstadt, 2019). Nếu các nhà nghiên cứu đang xử lý những tập dữ liệu đặc biệt nhạy cảm, cách tiếp cận tốt nhất có thể là tránh hoàn toàn các mô hình AI dựa trên đám mây và chọn các công cụ AI chạy cục bộ. Một số trường đại học và tổ chức nghiên cứu hiện đang phát triển các môi trường AI cô lập (sandbox) cho phép nhà nghiên cứu sử dụng công cụ AI mà không phải truyền dữ liệu đến các máy chủ bên ngoài (Masso et al., 2025).

Vượt ra ngoài các lựa chọn của từng nhà nghiên cứu, tính minh bạch là yếu tố thiết yếu khi công cụ AI được sử dụng trong nghiên cứu có sự tham gia của con người. Nếu AI hỗ trợ dưới bất kỳ hình thức xử lý hay phân tích dữ liệu nào, các văn bản chấp thuận có hiểu biết (informed consent) cần nêu rõ điều này. Một mẫu phiếu chấp thuận được cấu trúc phù hợp có thể ghi: \emph{``Nghiên cứu này sử dụng các công cụ hỗ trợ bởi AI để phân tích dữ liệu đã được ẩn danh. Không có thông tin định danh cá nhân nào được xử lý, và mọi thông tin sẽ được giữ bí mật. Mô hình AI được sử dụng không lưu giữ dữ liệu cũng như không đóng góp vào quá trình huấn luyện trong tương lai.''} Khi các chính sách của tổ chức tiếp tục phát triển, nhiều hội đồng đạo đức và Hội đồng Đánh giá Đạo đức Nghiên cứu (IRB) đang tăng cường giám sát vai trò của AI trong việc xử lý dữ liệu nghiên cứu. Một số hội đồng yêu cầu nhà nghiên cứu công bố rõ ràng liệu công cụ AI có được dùng trong phân tích dữ liệu, phiên âm hay bất kỳ khía cạnh nào khác của quy trình nghiên cứu hay không. Việc không tuân thủ các yêu cầu này có thể dẫn đến vi phạm về đạo đức hoặc quy định, đặc biệt đối với các dự án liên quan đến nhóm đối tượng dễ bị tổn thương, nghiên cứu mật hoặc hợp tác quốc tế chịu sự điều chỉnh của nhiều chế độ bảo vệ dữ liệu.

An ninh dữ liệu là một cân nhắc quan trọng khác bên cạnh việc ẩn danh hóa. Không giống phần mềm nghiên cứu truyền thống vốn cho phép người dùng kiểm soát việc lưu trữ và xử lý dữ liệu, nhiều công cụ AI vận hành trên hạ tầng đám mây độc quyền, nơi nhà cung cấp có thể truy cập thông tin. Điều này khiến việc bảo đảm thông tin nhạy cảm được bảo vệ khỏi truy cập hoặc lạm dụng trái phép trở nên khó khăn. Một giải pháp tiềm năng cho các nhà nghiên cứu làm việc với dữ liệu bảo mật là sử dụng các mô hình AI tự lưu trữ (self-hosted), tức các mô hình ngôn ngữ lớn chạy trên máy chủ cục bộ hoặc mạng nội bộ của tổ chức mà không truyền dữ liệu ra bên ngoài. Ví dụ gồm các mô hình mã nguồn mở như LLaMA, Falcon và Mistral, có thể được tinh chỉnh và triển khai trong các môi trường tính toán an toàn. Một số trường đại học đang đầu tư vào hạ tầng AI riêng, cho phép giảng viên sử dụng công cụ AI đồng thời bảo đảm dữ liệu không rời khỏi sự kiểm soát của tổ chức.

Các nhà nghiên cứu xử lý các nghiên cứu độc quyền, mật hoặc chưa công bố cũng nên lưu ý cách các công cụ AI quản lý việc lưu trữ dữ liệu tạm thời. Một số nền tảng lưu giữ dữ liệu phiên làm việc cho mục đích vận hành, ngay cả khi không dùng chúng để huấn luyện. Việc kiểm tra xem nhà cung cấp AI có lưu giữ nhật ký hoặc bản ghi hội thoại hay không là rất quan trọng trước khi nhập tài liệu nhạy cảm. Trong một số trường hợp, các thỏa thuận hợp đồng, chẳng hạn Thỏa thuận Xử lý Dữ liệu (DPA) theo GDPR, có thể cung cấp thêm các bảo đảm pháp lý rằng dữ liệu sẽ không bị lưu trữ hoặc chia sẻ với bên thứ ba.

Bằng cách thực hiện những biện pháp phòng ngừa này, rà soát các thiết lập quyền riêng tư của AI, tránh nhập dữ liệu có thể nhận dạng được hoặc dữ liệu chưa công bố, sử dụng ẩn danh hóa một cách cẩn trọng, lựa chọn các công cụ có cơ chế bảo vệ quyền riêng tư nghiêm ngặt, và bảo đảm tính minh bạch trong việc xin sự đồng ý có hiểu biết, các nhà nghiên cứu có thể giảm thiểu rủi ro trong khi vẫn tận dụng được những năng lực của AI. Do bản chất phát triển nhanh chóng của quản trị AI, việc cập nhật thường xuyên các chính sách của tổ chức, khung pháp lý và các thực hành tốt nhất sẽ là điều thiết yếu để sử dụng AI một cách có trách nhiệm trong nghiên cứu học thuật.

\subsection{Các cân nhắc về sở hữu trí tuệ trong việc sử dụng AI tạo sinh ở môi trường học thuật}

Việc tích hợp AI tạo sinh vào nghiên cứu và viết học thuật đặt ra những thách thức đáng kể về sở hữu trí tuệ (intellectual property, IP). Không giống các công cụ nghiên cứu truyền thống, AI tạo sinh có thể tạo ra những phần văn bản đáng kể, tóm tắt tài liệu hiện có, thậm chí hỗ trợ xây dựng cấu trúc lập luận. Những khả năng này làm nảy sinh các câu hỏi phức tạp về quyền tác giả, quyền sở hữu và việc ghi nhận nguồn, những vấn đề ngày càng trở nên cấp thiết khi các tổ chức học thuật và khuôn khổ pháp lý đang chật vật theo kịp tốc độ phát triển nhanh chóng của AI. Trong khi các cuộc tranh luận rộng hơn vẫn tiếp diễn về việc huấn luyện AI trên các tài liệu có bản quyền có cấu thành vi phạm hay không (Lemley, 2024), phần này tập trung vào các vấn đề sở hữu trí tuệ mà nhà nghiên cứu, giảng viên và sinh viên phải đối mặt khi sử dụng AI tạo sinh trong bối cảnh học thuật.

Một trong những mối quan ngại chính là liệu nhà nghiên cứu sử dụng AI tạo sinh có giữ toàn quyền sở hữu đối với văn bản được tạo ra hay không. Luật bản quyền nhìn chung giả định rằng chỉ tác giả là con người mới có thể yêu cầu quyền sở hữu trí tuệ, điều này làm vấn đề trở nên phức tạp khi AI đóng góp đáng kể vào việc tạo ra nội dung. Tại Hoa Kỳ, Văn phòng Bản quyền đã phán quyết rằng các tác phẩm hoàn toàn do AI tạo ra không đủ điều kiện được bảo hộ bản quyền (U.S. Copyright Office, 2023). Tuy nhiên, các tác phẩm có sự hỗ trợ của AI có thể được bảo hộ nếu có đủ sự kiểm soát và đóng góp sáng tạo của con người, nhưng việc xác định này được thực hiện theo từng trường hợp cụ thể (U.S. Copyright Office, 2025). Thách thức then chốt là xác định thế nào là đóng góp sáng tạo của con người. Nếu một nhà nghiên cứu soạn thảo bài báo với sự hỗ trợ của AI nhưng chỉnh sửa và trau chuốt nội dung một cách đáng kể, bảo đảm rằng các lập luận, phân tích và kết luận là của chính mình, thì điều này nhiều khả năng được coi là quyền tác giả của con người. Tuy nhiên, nếu AI đóng vai trò chi phối trong việc cấu trúc tác phẩm hoặc tổng hợp thông tin, thì mức độ mà tác phẩm vẫn được coi là ``nguyên gốc'' theo luật bản quyền trở nên mơ hồ hơn.

Vấn đề này đặc biệt đáng chú ý trong các dự án nghiên cứu hợp tác, nơi nhiều tác giả cùng đóng góp cho một bài báo. Nếu một nhà nghiên cứu phụ thuộc nhiều vào AI trong khi những người khác thì không, thì việc này nên được công bố như thế nào? Một số nhà xuất bản học thuật hiện yêu cầu tác giả nêu rõ liệu các công cụ AI có được dùng để soạn thảo bài báo hay không, trong khi những nhà xuất bản khác cấm rõ ràng việc ghi AI là đồng tác giả. Quy định thứ hai phản ánh nguyên tắc rằng quyền tác giả đồng nghĩa với trách nhiệm đối với nội dung, và vì AI không thể chịu trách nhiệm giải trình nên không thể được coi là tác giả (Lemley, 2024). Tuy nhiên, các chính sách này không giải quyết được câu hỏi rộng hơn về việc mức độ sử dụng AI bao nhiêu thì còn chấp nhận được trước khi vai trò tác giả của nhà nghiên cứu bị suy giảm. Khi các tổ chức và nhà xuất bản xây dựng những chính sách rõ ràng hơn, nhà nghiên cứu phải luôn cập nhật các hướng dẫn đang thay đổi và bảo đảm rằng những đóng góp của mình vẫn mang ý nghĩa trí tuệ.

Đây là lúc đáng nhắc đến khái niệm ``co-intelligence'' (trí tuệ cộng tác) của Ethan Mollick, khái niệm xem AI như một cộng tác viên chứ không phải sự thay thế cho quyền tác giả của con người (Mollick, 2024). Theo mô hình này, AI nâng cao hiệu quả và tính sáng tạo nhưng không đảm nhận vai trò tác giả chính. Cách tiếp cận này phù hợp với cách nhiều học giả hiện đang dùng AI, trau chuốt bản thảo, gợi ý các lập luận thay thế và cải thiện độ rõ ràng, thay vì tạo ra toàn bộ bài báo mà không có sự giám sát. Nếu AI được sử dụng trong khuôn khổ trí tuệ cộng tác này, tính minh bạch trở nên then chốt. Nhà nghiên cứu nên ghi lại việc AI đã tham gia vào công trình của mình như thế nào, phân biệt các đầu vào do AI tạo ra với những đóng góp trí tuệ của chính họ. Việc ghi chép rõ ràng có thể giúp ngăn ngừa tranh chấp về quyền tác giả, đặc biệt trong các dự án hợp tác, nơi những người đóng góp khác nhau có thể dùng AI ở các mức độ khác nhau.

Để bảo vệ các đóng góp trí tuệ của mình, nhà nghiên cứu nên lưu giữ các bản thảo cho thấy nội dung do AI tạo ra đã được chỉnh sửa ra sao. Bằng cách trau chuốt văn bản do AI tạo ra để phản ánh lập luận của chính mình và bảo đảm các luận điểm then chốt vẫn mang tính nguyên gốc, các học giả có thể giảm thiểu nguy cơ AI lấn át quyền tác giả của họ (Lemley, 2024). Các tổ chức và nhóm nghiên cứu cũng nên thiết lập các hướng dẫn nội bộ về việc sử dụng AI trong các dự án hợp tác, xác định khi nào và bằng cách nào các đóng góp của AI cần được công bố.

Ngoài văn bản, AI tạo sinh ngày càng được dùng để tạo hình, đồ thị và sơ đồ khái niệm cho các ấn phẩm học thuật. Mặc dù điều này mang lại cách thức hiệu quả và tiết kiệm chi phí để tạo nội dung trực quan, nó cũng kéo theo các rủi ro về bản quyền. Một số hình ảnh do AI tạo ra có thể giống với tài liệu có bản quyền hiện có, làm dấy lên lo ngại về vi phạm vô ý (U.S. Copyright Office, 2025). Nhà nghiên cứu cần kiểm tra cẩn thận để bảo đảm rằng hình ảnh do AI tạo ra không sao chép các tác phẩm được bảo hộ, đặc biệt khi dùng các công cụ tạo ảnh dựa trên những tập dữ liệu lớn về phương tiện truyền thông hiện có. Một số nhà xuất bản hiện yêu cầu tác giả xác nhận rằng mọi hình do AI tạo ra đều tuân thủ các tiêu chuẩn bản quyền. Vì nhiều công cụ AI đặt gánh nặng bảo đảm tính nguyên gốc lên người dùng, nhà nghiên cứu nên thực hiện các biện pháp để xác minh sự tuân thủ, chẳng hạn chạy hình ảnh do AI tạo ra qua tìm kiếm ảnh ngược hoặc sử dụng các công cụ xác minh bản quyền.

Để giảm thiểu rủi ro vi phạm bản quyền, nhà nghiên cứu nên lựa chọn công cụ AI một cách cẩn trọng. Các nền tảng được huấn luyện trên các tập dữ liệu truy cập mở có thể ít tạo ra nội dung vi phạm hơn. Ngoài ra, khi có thể, nhà nghiên cứu có thể dùng các công cụ AI cho phép họ huấn luyện mô hình trên bộ dữ liệu của riêng mình thay vì dựa vào các mô hình đa dụng có nguồn không rõ ràng (Lemley, 2024). Cách tiếp cận này làm giảm khả năng vi phạm bản quyền vô ý và giúp kiểm soát tốt hơn dữ liệu được dùng để tạo nội dung.

Ngoài bản quyền, việc sử dụng AI tạo sinh trong nghiên cứu còn làm dấy lên những lo ngại về bằng sáng chế và bí mật thương mại. Nhiều công cụ AI hoạt động trên các nền tảng điện toán đám mây, nghĩa là bất kỳ thông tin nào được nhập vào đều có thể bị lưu trữ, xử lý hoặc được dùng để huấn luyện mô hình trong tương lai. Nếu một nhà nghiên cứu nhập các phát hiện nghiên cứu mới, các giả thuyết hoặc dữ liệu thực nghiệm vào một hệ thống AI, thông tin này có thể bị truy cập hoặc vô tình bị lộ. Điều này có hệ quả nghiêm trọng đối với bằng sáng chế, bởi việc công bố công khai một sáng chế trước khi nộp đơn đăng ký bằng sáng chế có thể khiến sáng chế đó không còn đủ điều kiện được cấp bằng ở nhiều khu vực pháp lý (U.S. Copyright Office, 2025). Các lĩnh vực phụ thuộc nhiều vào bảo hộ sở hữu trí tuệ, chẳng hạn công nghệ sinh học, kỹ thuật và phát triển phần mềm, cần đặc biệt thận trọng khi sử dụng AI tạo sinh trong giai đoạn nghiên cứu ban đầu.

Để bảo vệ các nghiên cứu có thể được cấp bằng sáng chế, các học giả nên tránh nhập các phát hiện chưa công bố vào hệ thống AI, trừ khi họ chắc chắn rằng công cụ đó không lưu trữ hoặc sử dụng dữ liệu đầu vào của họ để huấn luyện. Hiện nay, một số nền tảng AI cung cấp các phiên bản dành cho doanh nghiệp hoặc học thuật, bảo đảm quyền riêng tư và an ninh dữ liệu. Khi có thể, các nhà nghiên cứu nên ưu tiên các phiên bản này hơn các mô hình miễn phí, công khai. Ngoài ra, các tổ chức nên xây dựng các thực hành tốt nhất cho nghiên cứu có hỗ trợ của AI, bảo đảm giảng viên và sinh viên hiểu rõ những rủi ro liên quan đến việc tiết lộ thông tin nhạy cảm cho các mô hình AI (U.S. Copyright Office, 2025).

Ngay cả khi quyền sở hữu trí tuệ không bị đe dọa, AI tạo sinh vẫn đặt ra những thách thức trong việc duy trì tính nguyên bản của công trình học thuật. Mặc dù luật bản quyền không bảo hộ bản thân các ý tưởng, nó bảo hộ cách thể hiện cụ thể của những ý tưởng đó. Nếu AI được dùng để tạo ra các bản tóm tắt, lập luận hoặc khung khái niệm, sẽ có nguy cơ nhiều người dùng cùng làm về các chủ đề tương tự nhận được những kết quả giống nhau. Vì AI tạo văn bản dựa trên các mẫu hình trong các tác phẩm hiện có, câu trả lời của nó có thể phản ánh cách diễn đạt quy ước, dẫn đến lo ngại về sự đồng nhất hóa của văn phong học thuật. Nếu nhiều nhà nghiên cứu dựa vào AI để tóm tắt tài liệu hoặc cấu trúc lập luận, diễn ngôn học thuật có thể trở nên đồng đều hơn, với ít phong cách viết riêng biệt hoặc những cách diễn giải mới mẻ hơn.

Các tổ chức, cơ quan tài trợ và nhà xuất bản đang bắt đầu đưa ra các chính sách rõ ràng hơn về việc sử dụng AI trong nghiên cứu học thuật. Một số trường đại học yêu cầu giảng viên khai báo liệu AI có được sử dụng trong các hồ sơ xin tài trợ, đề xuất nghiên cứu hoặc các ấn phẩm nộp để thẩm định nội bộ của tổ chức hay không. Tương tự, các cơ quan tài trợ có thể áp dụng các chính sách nhằm bảo đảm nghiên cứu có hỗ trợ của AI đáp ứng các chuẩn mực đạo đức và pháp lý. Các nhà nghiên cứu nên theo dõi các hướng dẫn đang không ngừng thay đổi này và bảo đảm rằng việc sử dụng AI của mình phù hợp với kỳ vọng của tổ chức và nghề nghiệp.

Khi các tòa án tiếp tục xem xét cách luật bản quyền áp dụng đối với nội dung do AI tạo ra, bối cảnh pháp lý có thể thay đổi. Cho đến khi có các quy định rõ ràng hơn, các học giả nên thận trọng khi tích hợp AI vào công việc của mình. Bằng cách ghi chép lại việc sử dụng AI, xác minh tính nguyên bản và bảo đảm tuân thủ các luật sở hữu trí tuệ, các nhà nghiên cứu có thể tối đa hóa lợi ích của AI đồng thời tránh được những cạm bẫy pháp lý tiềm ẩn.

\subsection{Những điểm cốt yếu về tuân thủ}

Việc tuân thủ khi sử dụng AI trong học thuật không chỉ dừng ở sở hữu trí tuệ và bảo vệ dữ liệu. Phần lớn phụ thuộc vào yêu cầu của tổ chức, hướng dẫn nộp bài của tạp chí và quy định về hồ sơ xin tài trợ. Không giống các khung pháp lý vốn tương đối ổn định, chính sách của các tổ chức về AI vẫn đang tiếp tục phát triển, nên các nhà nghiên cứu càng cần cập nhật những quy định cụ thể áp dụng cho công việc của mình.

Một trong những yếu tố đầu tiên cần cân nhắc là chính sách của trường đại học. Nhiều cơ sở học thuật đang trong quá trình xác định lập trường của mình đối với nghiên cứu và viết có hỗ trợ của AI. Một số trường hiện yêu cầu công bố rõ ràng khi AI được dùng để soạn thảo bài báo nghiên cứu, hồ sơ xin tài trợ hoặc tài liệu giảng dạy, trong khi những trường khác áp đặt hạn chế đối với nội dung do AI tạo ra, đặc biệt trong đánh giá sinh viên. Trước khi đưa AI vào công việc học thuật, các nhà nghiên cứu nên kiểm tra chính sách chính thức của cơ sở mình hoặc tham vấn các hội đồng khoa có liên quan. Trong một số trường hợp, việc không công bố việc sử dụng AI, nhất là trong các ứng dụng có mức độ rủi ro cao như hồ sơ xin tài trợ, có thể dẫn đến những vấn đề đạo đức hoặc rắc rối về hành chính.

Các cơ quan tài trợ cũng đang điều chỉnh quy định để đáp ứng vai trò ngày càng lớn của AI trong nghiên cứu. Một số đơn vị cấp kinh phí hiện yêu cầu người nộp đơn nêu rõ liệu AI có được dùng để soạn thảo đề xuất, thực hiện các phân tích sơ bộ hay rà soát tài liệu hay không. Những đơn vị khác có thể áp đặt giới hạn chặt chẽ hơn, đặc biệt đối với các dự án liên quan đến thông tin mật, dữ liệu nhạy cảm hoặc tuân thủ quy định. Do hồ sơ xin tài trợ phải chịu sự thẩm định nghiêm ngặt, việc không phù hợp với chính sách của một cơ quan có thể làm mất cơ hội nhận tài trợ. Trước khi nộp đề xuất, các nhà nghiên cứu nên xem xét kỹ hướng dẫn cụ thể của từng đơn vị tài trợ để bảo đảm việc sử dụng AI nằm trong phạm vi được phép.

Các nhà xuất bản học thuật là một yếu tố quan trọng khác. Nhiều tạp chí đang tích cực cập nhật chính sách biên tập để giải quyết vai trò của AI trong viết học thuật. Một số yêu cầu tác giả công bố liệu có sử dụng AI hỗ trợ soạn thảo hay không, trong khi những tạp chí khác cấm ghi AI là đồng tác giả, nhằm củng cố nguyên tắc rằng tác giả là con người phải chịu trách nhiệm về nội dung. Một số tạp chí có thể cấm hoàn toàn mọi việc sử dụng AI tạo sinh, hoặc hạn chế theo một cách khác. Ngoài ra, các tạp chí có thể quy định việc sử dụng hình minh họa, biểu đồ hoặc hình ảnh do AI tạo ra, đặc biệt nếu chúng làm dấy lên lo ngại về tính nguyên bản hoặc vi phạm bản quyền. Trước khi nộp bản thảo, tác giả nên xem xét chính sách của tạp chí để xác định liệu nội dung do AI tạo ra có phải được khai báo hay không và có áp dụng hạn chế cụ thể nào hay không. Bỏ qua các hướng dẫn này có thể dẫn đến chậm trễ, bị từ chối, cũng như các vấn đề đạo đức và thậm chí pháp lý.

Ngoài nghiên cứu và xuất bản, các chính sách về AI cũng đang xuất hiện trong các hội nghị học thuật và quy trình phản biện đồng cấp (peer review). Một số hội nghị đã bắt đầu quy định liệu các bản tóm tắt, đề xuất hoặc bài thuyết trình do AI tạo ra có được chấp nhận hay không. Mặc dù các công cụ AI có thể hữu ích trong việc trau chuốt lập luận và nâng cao độ rõ ràng, một số tổ chức học thuật đã bày tỏ lo ngại về việc lạm dụng AI trong các bài nộp. Tương tự, việc sử dụng AI trong phản biện đồng cấp vẫn còn gây tranh cãi: một số người phản biện dùng AI để tóm tắt bài báo, nhưng nhiều tổ chức học thuật không khuyến khích dùng nó để đánh giá các bài nộp, vì có thể đưa vào những thiên kiến hoặc sự đơn giản hóa quá mức. Hiểu những chuẩn mực mới nổi này là điều thiết yếu đối với các nhà nghiên cứu thường xuyên tham gia phản biện đồng cấp hoặc các hội nghị.

Để vượt qua những chính sách đang thay đổi này, các nhà nghiên cứu nên có cách tiếp cận chủ động. Trước khi dùng AI trong nghiên cứu, viết lách hoặc giảng dạy, họ nên kiểm tra mọi hướng dẫn của tổ chức, chính sách của tạp chí hoặc hạn chế của đơn vị tài trợ. Nếu chính sách chưa rõ ràng, việc tìm kiếm sự làm rõ từ trưởng khoa, phòng pháp chế hoặc biên tập viên tạp chí có thể ngăn ngừa các rắc rối về sau. Khi còn nghi ngờ, minh bạch là cách tiếp cận tốt nhất: công bố việc sử dụng AI khi được yêu cầu và duy trì sự giám sát của con người đối với các khía cạnh then chốt của công việc học thuật sẽ bảo đảm cả sự tuân thủ lẫn liêm chính nghiên cứu.

\section{Những cân nhắc và thách thức về đạo đức}

Các khía cạnh đạo đức của AI tạo sinh trong môi trường học thuật vượt ra ngoài việc tuân thủ pháp luật và chính sách của tổ chức, thách thức những chuẩn mực lâu đời về quyền tác giả, lao động trí tuệ và sự công bằng trong công việc học thuật (Jabotinsky \& Sarel, 2023). Nhiều học giả, đặc biệt là những người lần đầu tiếp xúc với các công cụ này, cảm thấy băn khoăn về mặt đạo đức; một số người thậm chí ví việc viết với sự hỗ trợ của AI như ``gian lận'' (Dwivedi et al., 2023). Phản ứng này là dễ hiểu. Các truyền thống học thuật nhấn mạnh nỗ lực trí tuệ của cá nhân, và các học giả được đào tạo để ghi nguồn cẩn thận khi xây dựng trên những ý tưởng có sẵn. Giờ đây, AI tạo sinh có thể hỗ trợ mọi việc, từ soạn câu văn đến gợi ý các khung lý thuyết, làm nảy sinh những câu hỏi căn bản về ranh giới của tư duy độc đáo và đóng góp trí tuệ.

Tuy nhiên, việc nhìn nhận việc sử dụng AI theo lối nhị phân, tức là có đạo đức hay phi đạo đức, là một sự đơn giản hóa quá mức. Những thách thức đạo đức cấp bách hơn không xuất phát từ bản thân việc sử dụng AI, mà từ các vấn đề gắn liền với cách thức hoạt động của các mô hình AI. Chúng bao gồm thiên kiến trong dữ liệu huấn luyện, tính thiếu minh bạch của các phản hồi do AI tạo ra, và nguy cơ phụ thuộc quá mức (Dwivedi et al., 2023). Những thách thức này đòi hỏi một cách tiếp cận tinh tế và có tính phản biện khi tích hợp AI vào công việc học thuật.

Một trong những mối quan ngại đạo đức quan trọng nhất là thiên kiến. Các mô hình AI được huấn luyện trên những tập dữ liệu khổng lồ gồm nội dung do con người tạo ra, nên không thể tránh khỏi việc thừa hưởng các thiên kiến có trong những nguồn đó (Jabotinsky \& Sarel, 2023). Do vậy, AI tạo sinh có thể củng cố những bất bình đẳng trong học thuật bằng cách trích dẫn không cân xứng các học giả nổi tiếng và các tổ chức danh tiếng, đồng thời gạt ra ngoài lề những nghiên cứu ít được biết đến. Vấn đề này đặc biệt đáng lo ngại trong môi trường học thuật, nơi sự đa dạng trí tuệ và việc tiếp cận nhiều quan điểm khác nhau có vai trò then chốt đối với sự tiến bộ của tri thức.

Chẳng hạn, khi được yêu cầu tóm tắt các nghiên cứu về một chủ đề nhất định, các công cụ AI có thể ưu tiên những bài báo được trích dẫn rộng rãi thay vì chọn ra những công trình đổi mới nhất hoặc có khả năng làm thay đổi lĩnh vực (Dwivedi et al., 2023). Xu hướng này củng cố các thứ bậc học thuật hiện có, khiến những tiếng nói thống trị càng thêm vững chắc, đồng thời gây khó khăn hơn cho các học giả mới nổi hoặc ít được công nhận về mặt thể chế trong việc được nhìn thấy. Vấn đề không chỉ dừng ở trích dẫn. Việc AI dựa vào các khuôn mẫu trong quá khứ khiến nó có xu hướng đưa ra những diễn giải theo lối quen thuộc thay vì tạo ra những hiểu biết mới. Điều này có thể dẫn đến sự đình trệ của diễn ngôn học thuật, khi các học giả dựa vào những bản tóm tắt do AI tạo ra vốn chỉ tái chế các quan điểm đang thịnh hành thay vì thách thức các mô hình tư duy hiện có.

Thiên kiến này có những hệ quả đặc biệt đáng kể đối với các nhà nghiên cứu giai đoạn đầu sự nghiệp và các học giả đến từ những khu vực hoặc ngành học ít được đại diện (Jabotinsky \& Sarel, 2023). Nếu các công cụ nghiên cứu có sự hỗ trợ của AI tiếp tục ưu tiên những công trình được trích dẫn nhiều từ các trung tâm học thuật thống trị, những tiếng nói thay thế sẽ càng khó được ghi nhận. Dù AI có thể thúc đẩy hiệu quả nghiên cứu, nó cũng có nguy cơ củng cố hiện trạng thay vì nuôi dưỡng sự đa dạng trí tuệ vốn thiết yếu cho tiến bộ học thuật. Để đối phó với điều này, các nhà nghiên cứu phải luôn cảnh giác trong việc bổ sung các tài liệu tham khảo do AI tạo ra bằng sự khám phá tài liệu của chính mình, bảo đảm rằng việc tiếp cận các nguồn học thuật của họ vẫn mang tính phản biện, bao trùm và toàn diện.

Ngoài thiên kiến trong nghiên cứu, sự thiếu minh bạch của AI đặt ra một thách thức đạo đức khác. Nhiều mô hình AI tạo sinh vận hành như những ``hộp đen'', nghĩa là người dùng không hiểu đầy đủ vì sao một mô hình AI đưa ra một phản hồi cụ thể, nó đang lấy thông tin từ những nguồn nào, hoặc liệu các yếu tố bên ngoài có ảnh hưởng đến đầu ra của nó hay không (Jabotinsky \& Sarel, 2023). Sự mờ đục này đặc biệt có vấn đề trong nghiên cứu học thuật, nơi tính minh bạch và khả năng tái lập là những nguyên tắc nền tảng. Các học giả phải có khả năng truy nguyên nguồn gốc các tài liệu của mình và biện minh cho kết luận của mình, nhưng nội dung do AI tạo ra lại thiếu cấu trúc trích dẫn rõ ràng và những lời giải thích về phương pháp mà các thực hành nghiên cứu truyền thống đòi hỏi.

Chẳng hạn, một số mô hình AI áp đặt các hạn chế về nội dung hoặc điều chỉnh phản hồi dựa trên những cơ chế bảo vệ được lập trình sẵn. Điều này có thể xuất phát từ động cơ chính trị hoặc ý thức hệ (Dwivedi et al., 2023). Một ví dụ nổi bật xảy ra với DeepSeek, một mô hình AI do Trung Quốc phát triển, vốn đưa ra những phản hồi mơ hồ hoặc bị kiểm duyệt khi được hỏi về các sự kiện lịch sử ở Trung Quốc. Những trường hợp như vậy cho thấy đầu ra của AI có thể bị định hình một cách tinh vi bởi các cân nhắc phi học thuật, khiến các học giả bắt buộc phải đánh giá một cách phản biện tài liệu do AI tạo ra thay vì chấp nhận nó theo giá trị bề mặt. Do các nhà nghiên cứu không phải lúc nào cũng biết những thiên kiến nào được gắn vào một mô hình, điều thiết yếu là tiếp cận các gợi ý nghiên cứu do AI tạo ra bằng sự thẩm xét kỹ lưỡng, đối chiếu các luận điểm với những nguồn đã được xác minh.

AI tạo sinh cũng làm nảy sinh những mối quan ngại đạo đức trong giảng dạy và đánh giá học thuật. Vai trò của AI trong giáo dục đang mở rộng, và dù nó có thể là một công cụ có giá trị cho việc học tập cá nhân hóa, tự động hóa chấm điểm và phát triển chương trình học, nó phải được sử dụng thận trọng để bảo vệ liêm chính học thuật. Một trong những mối lo cấp bách nhất là việc sinh viên ngày càng phụ thuộc vào các bài luận, bài tập và thậm chí cả bài thi do AI tạo ra (Dwivedi et al., 2023). Xu hướng này đặt các nhà giáo dục vào một thế lưỡng nan: nên cấm nghiêm ngặt nội dung do AI tạo ra, hay nên khuyến khích sinh viên sử dụng AI một cách có trách nhiệm, kèm theo các hướng dẫn rõ ràng về thời điểm và cách thức sử dụng phù hợp?

Một giải pháp khả dĩ nằm ở năng lực hiểu biết về AI có đạo đức, dạy sinh viên không chỉ cách sử dụng AI mà còn cách tiếp cận nó một cách phản biện (Jabotinsky \& Sarel, 2023). Thay vì cấm hoàn toàn, các trường đại học có thể yêu cầu sinh viên công khai sự hỗ trợ của AI trong bài làm, tương tự như việc trích dẫn nguồn trong nghiên cứu truyền thống. Cách tiếp cận này thừa nhận vai trò của AI như một công cụ hợp pháp, đồng thời ngăn chặn các hành vi gian dối. Tuy nhiên, các chính sách như vậy phải được xây dựng cẩn trọng để tránh tạo ra gánh nặng không đáng có cho sinh viên hoặc giảng viên, đặc biệt tại những cơ sở mà nguồn lực về hiểu biết AI chưa phổ biến.

Một mối lo ngại khác là kỹ năng tư duy phản biện và viết của sinh viên có thể bị xói mòn nếu họ quá phụ thuộc vào AI. Nếu sinh viên thường xuyên dựa vào AI để cấu trúc lập luận hoặc tạo ra các nhận định, họ có thể không phát triển được sự chặt chẽ trí tuệ cần thiết cho việc tham gia học thuật chuyên sâu (Dwivedi et al., 2023). Để giải quyết vấn đề này, giảng viên nên thiết kế các hình thức đánh giá đòi hỏi tư duy độc lập và sự tham gia chủ động, thi vấn đáp, bài tập viết tay và thảo luận trên lớp không dùng AI là những cách khả thi để bảo đảm sinh viên tiếp tục phát triển các kỹ năng phân tích thiết yếu.

Các trường đại học và nhà xuất bản học thuật đang bắt đầu ứng phó với những mối quan ngại đạo đức này bằng cách thiết lập các chính sách rõ ràng hơn về việc sử dụng AI (Jabotinsky \& Sarel, 2023). Như đã đề cập, một số tạp chí hiện yêu cầu tác giả công khai việc có sử dụng hỗ trợ viết bằng AI trong bài báo hay không, trong khi những tạp chí khác cấm ghi AI là đồng tác giả nhằm củng cố nguyên tắc rằng tác giả phải chịu trách nhiệm về công trình của mình. Tương tự, các cơ quan tài trợ ngày càng yêu cầu minh bạch về việc sử dụng AI trong hồ sơ xin tài trợ. Những chính sách mới nổi này cho thấy sự thừa nhận ngày càng tăng ở cấp độ thể chế rằng AI đang định hình lại công việc học thuật, và các nhà nghiên cứu phải thích ứng với các chuẩn mực đang thay đổi.

Rốt cuộc, AI tạo sinh không phải là một lợi ích tuyệt đối cũng không phải là mối đe dọa vốn có đối với liêm chính học thuật. Tác động của nó phụ thuộc vào cách thức sử dụng. Nếu giới học giả luôn tham gia một cách phản biện, nhận thức rõ những hạn chế của AI, đặt câu hỏi về các thiên kiến của nó, và bổ sung kết quả đầu ra của nó bằng nghiên cứu nghiêm cẩn, thì những công cụ này có thể đóng vai trò là trợ lý hữu ích thay vì là nguồn gốc của các vấn đề đạo đức. Tuy nhiên, nếu không giải quyết những thách thức này, có thể dẫn đến một bối cảnh học thuật trong đó AI không chỉ đẩy nhanh nghiên cứu mà còn làm méo mó nó, củng cố các bất bình đẳng sẵn có và làm suy yếu các nguyên tắc nền tảng của hoạt động nghiên cứu học thuật. Bằng cách tiếp cận AI với sự hoài nghi, nhận thức đạo đức và sự nhấn mạnh vào chuyên môn của con người, giới học thuật có thể khai thác tiềm năng của AI đồng thời bảo vệ những giá trị làm nên tiến bộ trí tuệ.

AI trong đánh giá học thuật: Thách thức và thực hành tốt nhất cho phản biện đồng cấp (peer review), đánh giá tài trợ và đánh giá sinh viên

Việc tích hợp AI vào các quy trình đánh giá học thuật, từ phản biện đồng cấp (peer review) và đánh giá tài trợ đến chấm điểm sinh viên, đặt ra những câu hỏi phức tạp về tính công bằng, thiên kiến, tính minh bạch và trách nhiệm giải trình. Dù AI có thể tinh giản các quy trình này, nâng cao hiệu quả và cung cấp những hiểu biết giá trị, nó cũng đưa vào những rủi ro cần được quản lý cẩn trọng để duy trì liêm chính học thuật.

Vai trò của AI trong phản biện đồng cấp (peer review) đặc biệt gây tranh cãi. Dù phản biện đồng cấp (peer review) vẫn là nền tảng của xuất bản học thuật, đây cũng là một quy trình không hoàn hảo và tốn thời gian, chịu ảnh hưởng của các thiên kiến và sự thiếu nhất quán của con người. Các công cụ hỗ trợ bởi AI đã được đề xuất như một cách để giảm bớt một phần những gánh nặng này. Chúng có thể tóm tắt bản thảo, làm nổi bật các lập luận chính, xác định các sai sót phương pháp luận và gắn cờ các mối quan ngại đạo đức tiềm tàng, giúp người phản biện đánh giá các bài nộp hiệu quả hơn. Một số công cụ thậm chí so sánh bản thảo mới với các tài liệu hiện có để phát hiện tính mới, sự trùng lặp hoặc khả năng đạo văn. Tuy nhiên, khả năng thực hiện các nhiệm vụ này một cách khách quan của AI còn hạn chế. Nó vận hành dựa trên các mẫu thống kê và xác suất ngôn ngữ thay vì sự thấu hiểu sâu sắc, khiến nó dễ bỏ sót những đóng góp lý thuyết tinh tế hoặc các cải tiến phương pháp luận.

Hơn nữa, các công cụ phản biện hỗ trợ bởi AI có thể củng cố các thứ bậc học thuật hiện có thay vì thúc đẩy đánh giá công bằng. Vì các mô hình AI tạo sinh được huấn luyện trên các sản phẩm học thuật trong quá khứ, chúng có xu hướng phản chiếu các mẫu trích dẫn đã được thiết lập và các thiên kiến thể chế, ưu ái nghiên cứu từ các trường đại học danh tiếng hoặc tạp chí có tác động cao trong khi bỏ qua những đóng góp ít được biết đến hơn nhưng có giá trị không kém. Điều này làm dấy lên mối lo ngại về tính công bằng và đa dạng trong xuất bản học thuật. Nếu AI được đưa vào phản biện đồng cấp (peer review), các tạp chí và nhà xuất bản phải bảo đảm rằng việc sử dụng nó không vô tình củng cố các thiên kiến mang tính cấu trúc. Tính minh bạch là then chốt: nếu AI được dùng để tạo bản tóm tắt phản biện hoặc các khuyến nghị, người phản biện cần biết các công cụ này ảnh hưởng đến đánh giá của họ như thế nào. AI không nên thay thế nhiệm vụ thiết yếu là phán đoán học thuật mang tính phản biện, cũng không nên được dùng như một thẩm quyền không thể nghi ngờ trong việc đánh giá bản thảo.

Một loạt thách thức tương tự nảy sinh khi AI được sử dụng trong đánh giá đề xuất tài trợ. Các cơ quan tài trợ đang bắt đầu khám phá vai trò của AI trong việc hợp lý hóa quy trình thẩm định các đề xuất nghiên cứu. AI có thể được dùng để kiểm tra mức độ tuân thủ các hướng dẫn tài trợ, nhận diện xung đột lợi ích tiềm ẩn và so sánh các đề xuất với những dự án đã được tài trợ. Về lý thuyết, điều này có thể giúp việc đánh giá đề xuất tài trợ trở nên có hệ thống hơn và ít chịu ảnh hưởng bởi thiên kiến cá nhân hơn. Tuy nhiên, cũng như với phản biện đồng cấp (peer review), các công cụ đánh giá dựa trên AI có thể làm nảy sinh những mối lo ngại mới. Nếu được huấn luyện trên các quyết định tài trợ trong quá khứ, các mô hình AI có thể tái tạo những thiên kiến cũ, củng cố các chênh lệch hiện có trong phân bổ kinh phí nghiên cứu. Một số ngành, phương pháp luận hoặc tổ chức có thể được ưu ái một cách có hệ thống so với những ngành, phương pháp luận hoặc tổ chức khác, từ đó hạn chế cơ hội cho các nghiên cứu đổi mới hoặc phi truyền thống. Hơn nữa, các khía cạnh định tính của việc đánh giá đề xuất tài trợ, chẳng hạn như thẩm định tác động tiềm năng, tính độc đáo và tính khả thi của một dự án, đòi hỏi sự phán đoán của con người mà AI không thể tái tạo đầy đủ. Để bảo đảm công bằng, các cơ quan tài trợ nên coi các đánh giá do AI tạo ra mang tính tư vấn thay vì mang tính quyết định, cho phép người nộp đề xuất phản biện các đánh giá khi cần thiết và duy trì sự giám sát của con người trong các quyết định tài trợ cuối cùng.

Ngoài xuất bản học thuật và phân bổ kinh phí tài trợ, AI ngày càng được sử dụng trong đánh giá sinh viên và chấm thi. Các công cụ AI có thể hỗ trợ giảng viên chấm bài luận, đánh giá bài thi và đưa ra phản hồi về bài làm của sinh viên. Đối với các bài đánh giá có cấu trúc, chẳng hạn như bài kiểm tra trắc nghiệm, việc chấm điểm dựa trên AI có thể mang lại hiệu quả đáng kể. Đối với các bài tập mở, các công cụ AI có thể phân tích ngữ pháp, tính mạch lạc và thậm chí cả cấu trúc lập luận, đưa ra những đánh giá sơ bộ để giảng viên tinh chỉnh thêm. AI cũng có tiềm năng trong đánh giá hình thành, khi nó có thể tạo ra phản hồi tức thì cho sinh viên, giúp họ nhận ra những điểm cần cải thiện trước khi nộp bài cuối cùng.

Tuy nhiên, không thể bỏ qua những hạn chế của AI trong đánh giá. AI không thể đánh giá đầy đủ tính độc đáo, chiều sâu phân tích hay tư duy phản biện, những phẩm chất cốt lõi của giáo dục đại học. Nếu bị lạm dụng, các công cụ chấm điểm bằng AI có nguy cơ biến việc đánh giá sinh viên thành nhận dạng mẫu ở bề mặt thay vì sự tiếp cận thực sự với bài làm của sinh viên. Hơn nữa, việc AI dựa vào dữ liệu trong quá khứ đồng nghĩa với việc nó có thể củng cố các cấu trúc diễn ngôn thống trị, phạt những lối tư duy phi truyền thống hoặc các cách tiếp cận sáng tạo không phù hợp với các khuôn mẫu được kỳ vọng. Trong các ngành đề cao lập luận và sắc thái diễn giải, khả năng chấm điểm của AI vẫn còn đáng nghi ngờ, nói một cách nhẹ nhất.

Vai trò của AI trong việc xây dựng đề thi cũng đặt ra những cân nhắc sư phạm quan trọng. AI có thể tạo ra các câu hỏi trắc nghiệm, đề bài luận và thậm chí cả các tình huống nghiên cứu điển hình dựa trên tài liệu môn học. Dù đây có thể là điểm khởi đầu hữu ích, giảng viên phải bảo đảm rằng các bài đánh giá do AI tạo ra phù hợp với mục tiêu học tập và không chỉ kiểm tra khả năng ghi nhớ thông tin. AI cũng có thể hỗ trợ phân tích các đề thi trước đây để xác định những nội dung sinh viên gặp khó khăn, giúp giảng viên tinh chỉnh chiến lược giảng dạy. Tuy nhiên, việc dựa vào các đề thi do AI tạo ra mà không có sự rà soát cẩn thận của con người có nguy cơ tạo ra những bài đánh giá hời hợt hoặc lặp lại, không đo lường được các kết quả học tập sâu hơn.

Liêm chính học thuật là một mối quan ngại lớn khác. Khi các công cụ AI ngày càng tinh vi, sinh viên có ngày càng nhiều khả năng tiếp cận sự trợ giúp do AI tạo ra cho bài tập và bài thi. Điều này đặt ra thách thức cho các giảng viên muốn bảo đảm thực hành đánh giá công bằng. Các công cụ phát hiện đạo văn truyền thống không phải lúc nào cũng hiệu quả trong việc nhận diện nội dung do AI tạo ra, làm dấy lên câu hỏi về cách các tổ chức nên xử lý công việc học thuật có sự hỗ trợ của AI. Một số trường đại học đang bắt đầu yêu cầu sinh viên khai báo việc sử dụng AI trong bài tập, tương tự như việc trích dẫn nguồn trong nghiên cứu thông thường. Dù cách tiếp cận này thúc đẩy tính minh bạch, nó cũng đặt ra câu hỏi về mức độ hỗ trợ của AI nào là chấp nhận được. Nếu AI giúp trau chuốt cấu trúc câu nhưng không tạo ra nội dung thực chất, liệu có cần khai báo không? Nếu AI gợi ý toàn bộ khung lập luận, liệu điều đó có cấu thành một lợi thế không công bằng? Những vấn đề này nhiều khả năng sẽ đòi hỏi các chính sách học thuật mới và sự thay đổi trong cách đánh giá tính độc đáo ở bài làm của sinh viên.

Thay vì cấm đoán hoàn toàn việc sử dụng AI, các tổ chức có thể được lợi khi thúc đẩy năng lực hiểu biết về AI có đạo đức trong sinh viên. Dạy sinh viên cách sử dụng AI như một công cụ bổ trợ, thay vì thay thế cho tư duy độc lập, có thể giúp ngăn ngừa hành vi sai phạm học thuật, đồng thời vẫn cho phép sinh viên hưởng lợi từ năng lực của AI. Giảng viên nên thiết kế các bài đánh giá nhấn mạnh sự tham gia mang tính phân tích, chẳng hạn như bài viết trên lớp, thi vấn đáp và các bài tập đòi hỏi viết nháp lặp đi lặp lại kèm phản hồi. Những cách tiếp cận này có thể giúp bảo đảm sinh viên phát triển kỹ năng tư duy phản biện và viết, thay vì phụ thuộc thụ động vào nội dung do AI tạo ra.

Khi AI ngày càng được tích hợp vào phản biện đồng cấp (peer review), đánh giá đề xuất tài trợ và đánh giá sinh viên, các tổ chức phải thiết lập những hướng dẫn rõ ràng để bảo đảm việc sử dụng AI củng cố thay vì làm suy yếu liêm chính học thuật. Tính minh bạch là then chốt: học giả, sinh viên và người phản biện phải hiểu AI đang được sử dụng như thế nào và những giới hạn của nó nằm ở đâu. Dù AI có thể hỗ trợ các quy trình này, nó không nên thay thế phán đoán của con người, và các đầu ra của nó cũng không nên được coi là những đánh giá dứt khoát về chất lượng học thuật. Bằng cách duy trì sự giám sát, nuôi dưỡng năng lực hiểu biết về AI và thiết kế các phương pháp đánh giá ưu tiên sự tham gia mang tính phản biện, giới học thuật có thể tận dụng lợi ích của AI trong khi vẫn giữ vững các chuẩn mực về công bằng, tính độc đáo và sự nghiêm cẩn học thuật.

Kết luận

Việc tích hợp AI tạo sinh vào môi trường học thuật vừa mang lại cơ hội vừa đặt ra thách thức, đòi hỏi các học giả, các tổ chức và các nhà hoạch định chính sách phải định hướng trong một bối cảnh pháp lý, đạo đức và thực tiễn đang không ngừng biến đổi. Khi các công cụ AI ngày càng tinh vi, chúng sẽ định hình ngày càng sâu sắc cách thức nghiên cứu được tiến hành, cách văn bản học thuật được tạo ra, cách sinh viên học tập và cách công trình học thuật được đánh giá. Tuy nhiên, những tiến bộ này cũng kéo theo các rủi ro cần được giải quyết, những bất định về sở hữu trí tuệ, các mối lo ngại về quyền riêng tư dữ liệu, các thiên kiến ẩn trong nội dung do AI tạo ra, và nguy cơ xói mòn tư duy phản biện cùng tính độc đáo trong công việc học thuật.

Các khung quy định, chính sách của tổ chức và chuẩn mực nghề nghiệp vẫn đang thích nghi với sự phát triển nhanh chóng của AI. Dù các đạo luật hiện hành về sở hữu trí tuệ, bảo vệ dữ liệu và liêm chính học thuật có đưa ra một số định hướng, chúng vẫn chưa phản ánh đầy đủ những phức tạp do AI tạo sinh đặt ra. Các nhà nghiên cứu và nhà giáo dục cần chủ động tìm hiểu cách các yêu cầu pháp lý và thể chế đang thay đổi này áp dụng vào công việc của mình, bảo đảm tuân thủ các luật về quyền riêng tư, các quy định bảo hộ sở hữu trí tuệ và các chính sách công bố thông tin của tổ chức. Tính minh bạch sẽ là yếu tố then chốt, đặc biệt trong những bối cảnh mà AI đóng vai trò đáng kể trong việc soạn thảo, phân tích hoặc ra quyết định.

Vượt ra ngoài việc tuân thủ pháp luật, các cân nhắc đạo đức phải luôn giữ vị trí trung tâm trong việc sử dụng AI ở môi trường học thuật. Những vấn đề như thiên kiến trong nội dung do AI tạo ra, tính thiếu minh bạch của quá trình ra quyết định của AI, và tính công bằng của việc đánh giá sinh viên với sự hỗ trợ của AI đều đòi hỏi sự xem xét một cách phản biện. Dù AI có thể đẩy nhanh nghiên cứu, nâng cao hiệu quả và hỗ trợ phân tích có giá trị, nó không được trở thành sự thay thế cho phán đoán, óc sáng tạo hay sự nghiêm cẩn trí tuệ của con người. Các học giả phải đánh giá một cách phản biện các kết quả đầu ra do AI tạo ra, đối chiếu chéo các nguồn, và bảo đảm rằng AI được dùng như một công cụ để nâng cao, chứ không thay thế, chuyên môn học thuật.

Đối với các tổ chức, thách thức sẽ là xây dựng những chính sách cân bằng giữa đổi mới và liêm chính học thuật. Các trường đại học phải xây dựng những hướng dẫn rõ ràng về vai trò của AI trong nghiên cứu, giảng dạy và đánh giá, bảo đảm rằng hiểu biết về AI (AI literacy) trở thành một phần cốt lõi của đào tạo học thuật. Tương tự, các nhà xuất bản học thuật và các cơ quan tài trợ phải hoàn thiện chính sách của mình để đáp ứng vai trò ngày càng lớn của AI trong công trình học thuật, đồng thời duy trì các tiêu chuẩn cao về quyền tác giả, ghi nhận nguồn và đạo đức nghiên cứu.

Rốt cuộc, AI tạo sinh không phải là một mối đe dọa hiện sinh, cũng không phải một lợi ích được bảo đảm đối với học thuật, đó là một công cụ mạnh mẽ mà tác động của nó phụ thuộc vào cách thức được tích hợp. Bằng cách tiếp cận AI với nhận thức phản biện, trách nhiệm đạo đức và cam kết đối với liêm chính học thuật, các học giả có thể khai thác tiềm năng của nó đồng thời bảo vệ những giá trị làm nền tảng cho hoạt động tìm hiểu học thuật. Tương lai của AI trong môi trường học thuật sẽ không chỉ được định hình bởi những tiến bộ công nghệ, mà còn bởi những lựa chọn của các tổ chức và cá nhân trong việc áp dụng, điều chỉnh và hoàn thiện cách sử dụng nó.
